\chapter{One document, from draft to decision}
\label{ch:worked-case}

One short research note makes the design concrete. An author, a parser, and a
reader each do something different with it. The note is a teaching case
assembled from the failure patterns in the
project's reader records; it is not a claim about a real model's performance or
a disguised forecast. Its purpose is to keep the nouns, numbers, and choices
visible while the later chapters generalize the method.

The note compares a baseline model with a constrained reconstruction and asks
whether the team should approve a larger comparative run. It contains the
things that make a technical document worth protecting: an equation, a table,
a citation, a numerical result, and a qualification about what the result does
not establish. It also contains the things that make a document hard to read:
private phase names, an implicit mechanism, a result separated from the
question it answers, and a request that appears only after several pages of
process description.

Figure~\ref{fig:worked-case-flow} is the route through the chapter. The return
arrow is bounded: a failed pre-human check goes back to a located authoring
decision before release. The final reader task is not the system's debugging
loop, and it never triggers an unending series of automatic rewrites.

\hvFigureWorkedCase

\section{The document before the tools}
\label{sec:case-before}

Imagine receiving the note as a PDF with no repository history. The first page
opens with a sentence such as ``Phase 2B evaluates the post-rescue path under
the frozen interface.'' The phrase is meaningful to the project team. A new
reader does not know what was rescued, which interface is frozen, or why the
comparison matters. The next page defines a loss function and presents a table
of parameter values. The decision request arrives near the end: ``If the
current branch clears the remaining checks, proceed to the expanded run.''

Nothing in those sentences is necessarily false. The problem is that the
reader has to infer the object of the comparison before deciding whether the
equation and table are relevant. In the internal records, this was the point
at which a mathematically serious draft could be rejected within minutes. The
reader did not ask for more adjectives. The reader asked what the team was
trying to change and what action the document supported.

\begin{table}[H]
\centering
\small
\begin{tabularx}{\textwidth}{@{}p{0.23\textwidth}L p{0.27\textwidth}@{}}
\toprule
Element in the source & The author's hidden assumption & Question the
reader actually needs answered \\
\midrule
Private phase label & The history of the project and the meaning of ``rescue'' &
What object is being compared, and why now? \\
Loss function & The target has already been defined elsewhere & What decision
does this quantity inform, and what are its units? \\
Parameter table & The rows are self-explanatory & Which rows are inputs, which
are estimates, and which are illustrative? \\
Qualification & The reader will connect it to the result two pages earlier &
What inference is supported, and where does it stop? \\
Final request & The implementation team knows the next phase & What exactly is
being approved, at what cost, and with what stopping rule? \\
\bottomrule
\end{tabularx}
\caption{The source contains the necessary material, but the reader must
reconstruct its order and purpose.}
\label{tab:case-before}
\end{table}

The distinction between an absent fact and a misplaced fact matters for the
design. A general language model can often make the first sentence smoother.
It cannot know whether ``post-rescue'' is a technical term, a project label,
or an accidental remnant unless the author supplies that meaning. A linter can
flag a long sentence. It cannot decide whether the loss function is the right
quantity for the funding request. The first review packet must therefore
show the source passage and ask a question that belongs to the author or the
reader.

\section{The first preflight and the first repair}
\label{sec:case-first-repair}

The proposed workflow starts with a preflight run, before any expensive reader
is asked to reconstruct the project. The system checks the brief and blueprint,
then tests the first unit against four questions: What decision is at stake?
What is the proposed mechanism? What evidence would change the decision? What
action is requested? These are not a score for prose quality. They are a
construction test for whether the document has delivered its economic purpose.

Suppose the preflight critic cannot recover the first and fourth answers. The
author opens a located finding. It includes the unit, its source location, the
reason it was selected, the missing dependency or evidence, and an editorial
question. The reason is descriptive rather than moral: the paragraph uses a
private phase label before introducing the research question and leaves the
action implicit. The author can accept the finding, reject it with a reason, or
rewrite the passage. A model suggestion may be displayed, but it is an option,
not the recorded decision. The unit remains inside the repair loop until the
preflight either passes or stops with an explicit unresolved choice.

\begin{repairpair}
\before{Phase 2B evaluates the post-rescue path under the frozen interface.}
\after{The team must decide whether the constrained reconstruction is ready
for a larger comparison. We compare it with the baseline model under the same
data and evaluation rule.}
\end{repairpair}

The revision does three jobs. It names the decision, identifies the two sides
of the comparison, and states the common condition that makes the comparison
interpretable. It does not claim that the constrained model is better. That
claim belongs after the evidence. It also does not delete the project history;
the history remains useful in the internal record, where it can be found by a
team member who needs it.

The second repair concerns an equation whose meaning was hidden by notation.
The source wrote
\begin{equation}
  L(\theta) = \sum_{t=1}^{T}
  \bigl(y_t-f_\theta(x_t)\bigr)^2,
  \label{eq:case-source-loss}
\end{equation}
without saying whether $L$ was a training objective, an evaluation statistic,
or an accounting identity. The equation is elementary; its role is not. The
author's revision introduces the units and purpose:

\begin{quote}
For the feasibility comparison, $y_t$ is the held-out outcome, $x_t$ is the
declared input record, and $f_\theta$ is the fitted reconstruction. We use
$L(\theta)$ only as a held-out prediction loss, measured in squared outcome
units. It is not a measure of reader comprehension and it does not determine
whether the team should approve deployment.
\end{quote}

The display itself remains unchanged. The surrounding prose does the work that
notation alone cannot do. A reader who cares about the model can now interpret
the terms; a decision-maker can see why the number is relevant but insufficient.

\section{Reconstructing the argument}
\label{sec:case-argument}

Once the first repairs are accepted, the author reorganizes the note around
five questions. The order is not a universal template; it fits this decision
because each answer supplies the premise for the next one.

\begin{enumerate}
\item \textbf{Decision.} The team must decide whether to fund a larger
      comparative run.
\item \textbf{Object.} The comparison is between a baseline and a constrained
      reconstruction applied to the same declared record.
\item \textbf{Mechanism.} The reconstruction changes the specified mapping
      from inputs to predictions while holding the evaluation rule fixed.
\item \textbf{Evidence.} A held-out loss, a protected-object comparison, and
      a reader test provide different pieces of evidence.
\item \textbf{Action.} The sponsor approves the next run only if the
      comparison is reproducible, the meaning-bearing objects are preserved,
      and the reader can recover the request.
\end{enumerate}

\begin{table}[H]
\centering
\small
\begin{tabularx}{\textwidth}{@{}p{0.18\textwidth}L p{0.26\textwidth}@{}}
\toprule
Argument step & Passage the revised note must contain & Observation that can
overturn it \\
\midrule
Decision & A one-sentence funding request with a budget and stopping
rule & The sponsor cannot identify a decision owner or a bounded action \\
Object & A definition of the baseline, reconstruction, population, and
evaluation window & The two versions use different records or targets \\
Mechanism & A short derivation or diagram showing what changes and what remains
fixed & The observed improvement comes from a changed data or scoring rule \\
Evidence & The held-out result, source location, and reader observation are
reported separately & The result is not reproducible or the reader cannot use it \\
Action & A proceed, revise, or stop recommendation tied to the observations &
The recommendation depends on an unstated preference \\
\bottomrule
\end{tabularx}
\caption{The argument becomes a sequence of answerable questions rather than a
sequence of project phases.}
\label{tab:case-argument}
\end{table}

The table is useful during revision, but it should not appear as a five-row
form in every document. In the finished note, the five answers can be ordinary
prose, a figure, or a short table wherever that form best carries the argument.
Humanvoice keeps the intermediate record because it lets an author explain why
a paragraph moved or why a table was retained. The reader sees the resulting
argument, not the form.

The revised note also changes the order of its numerical material. First it
states the decision and the target. Then it gives the equation and explains its
role. The table follows with a sentence that identifies inputs, estimates, and
illustrative values. Finally, the qualification appears beside the result it
qualifies. This is a small editorial change, but it reduces the number of
working assumptions the reader must hold at once.

\section{The objects the system protects}
\label{sec:case-protected}

The revision is not safe merely because the new prose sounds plausible. The
source contains objects whose alteration can change the scientific meaning.
For this case, the protected set is
\begin{equation}
  \mathcal P = \{\text{equations},\ \text{numbers},\ \text{labels},\
  \text{citations},\ \text{tables},\ \text{declared qualifications}\}.
  \label{eq:case-protected-set}
\end{equation}
Each member receives a source span, a normalized representation, and a
before-and-after comparison. The comparison is not a demand that every object
remain identical. It is a demand that a change become visible and receive an
owner.

For an object $p\in\mathcal P$, write $\Delta p=0$ when the normalized object
is unchanged, $\Delta p=\mathrm{explained}$ when the author has recorded a
meaningful change, and $\Delta p=\mathrm{unresolved}$ when the comparison
cannot establish correspondence. The release rule for the case is
\begin{equation}
  \forall p\in\mathcal P,\qquad
  \Delta p \ne \mathrm{unresolved}.
  \label{eq:case-protection-rule}
\end{equation}
This is a safety rule for the workflow, not a quality score. It does not say
that an unchanged equation is correct or that an explained change is wise. It
says that the system must not hide a broken correspondence behind a fluent
rewrite.

The note contains one numerical wedge, for example a change from $0.42$ to
$0.47$ in a displayed comparison. A character-level diff will find the change,
but it will not tell the author whether the value is a typo, a recalculation,
or a deliberate change in the sample. The protected record links the number to
its table row, unit, source calculation, and explanation. If the author cannot
provide that link, the revision returns to the author regardless of how much
the surrounding paragraph improved.

Citation protection has the same shape. The system resolves the key to a
bibliographic identity and records the sentence in which it appears. It does
not infer that the source supports the sentence. That second judgment remains
with the author and, for a load-bearing claim, an independent reviewer. A DOI
resolver can prevent a misspelled reference from surviving; it cannot turn a
nearby paper into evidence.

\section{The author's revision}
\label{sec:case-author-loop}

After the first pass, the author sees a queue of findings ordered by expected
reader consequence, not by the model's confidence. Three findings may appear:

\begin{description}[style=nextline,leftmargin=34mm,labelindent=0mm]
\item[Meaning] The paragraph names a result but does not say which decision it
supports.
\item[Register] A private phase name and two file paths make the opening
unintelligible outside the project.
\item[Surface] Three adjacent sentences repeat the same grammatical shape.
\end{description}

The author chooses the meaning finding first. This choice is important because
surface repair can make a weak paragraph more pleasant without making it more
useful. The author then decides whether the phase name has a legitimate public
meaning. If it does, it is defined once. If it does not, it is removed from the
reader-facing version and retained in the internal source map. Finally, the
author may vary the sentence rhythm, but only after the proposition is clear.

The queue records the decision in ordinary language: ``accepted: the reader
needs the action before the phase history'' is more useful than a status code.
The machine can store a code for searching, but the code is not the explanation
that another author or reviewer should have to read.

\begin{cardexample}{a rejected suggestion}
The editor proposes replacing the qualification ``the comparison does not
establish transport to a new population'' with ``the result is promising but
further work is needed.'' The second sentence is smoother and less useful. It
removes the population condition and gives no account of what further work
would answer the question. The author rejects the suggestion and rewrites the
qualification in terms of the missing population and the planned comparison.
\end{cardexample}

This example illustrates why a human decision belongs in the path. The issue
is not whether the proposed sentence is grammatical. It is whether the sentence
preserves the inference boundary that a reader could reasonably misunderstand.
An automatic system can point to the boundary, show the alternative, and
record the protected comparison. It cannot own the scientific judgment.

\section{The reader's turn}
\label{sec:case-reader}

After the preflight gates pass, the revised note is rendered and sent to a reader who has not seen the source
history, the finding queue, or the author's rejected suggestions. The reader
receives a short context statement: ``You are deciding whether to approve a
larger comparative run. Read the note as you would a proposal from a colleague
and answer the four questions below.'' The reader may annotate the document,
but the primary observations are timed and forced-choice so that a vague
impression becomes a usable record.

For reader $j$ and version $v$, let $A_{jv}=1$ when the reader identifies the
requested action, the comparison, and the principal qualification within the
declared time budget. Let $T_{jv}$ be reading and clarification time, and let
$R_v$ be the number of author feedback rounds before acceptance. The simple
case-level outcome is the vector
\begin{equation}
  O_v = \left(\frac{1}{J}\sum_{j=1}^{J}A_{jv},\quad
              \frac{1}{J}\sum_{j=1}^{J}T_{jv},\quad R_v,\quad
              M_v\right),
  \label{eq:case-reader-outcome}
\end{equation}
where $M_v$ counts unresolved or unauthorized changes to protected objects.
The vector is deliberately not collapsed into one score. A faster reader who
misses the qualification is not an improvement.

The first reader may answer ``revise'' even when the revised prose is fluent.
That is not a failed product demonstration. It is information about which
finding remains unresolved. The author receives the reader's question and
returns to the source, rather than asking the reader to explain the entire
document. A second reading then distinguishes a local repair from a change in
the argument. The process ends when the reader can make the declared decision
or when the sponsor invokes the stop rule.

\section{The economic reading of the case}
\label{sec:case-economics}

The case also makes the investment mechanism concrete. Suppose the incumbent
version takes the reader $T_0=3.0$ hours and requires $R_0=2.4$ feedback rounds
on average. A revised version takes $T_1=2.4$ hours and $R_1=1.8$ rounds. For
$N=100$ documents per year and a loaded reader cost of $c=150$ per hour, the
direct attention saving is
\begin{equation}
  B = N(T_0-T_1)c = 100(0.6)(150)=9{,}000.
  \label{eq:case-benefit}
\end{equation}
The values are teaching numbers. They show why the reader's time must be
measured directly and why a lower round count is not enough. If the revised
version adds a single hour of clarification for each document, the apparent
saving disappears. If $M_1>0$, the financial calculation is beside the point
until the protected change is explained.

The case also separates costs that are easy to hide. The author spends time
answering findings; the reader spends time reading and clarifying; a technical
owner spends time maintaining parsers and fixtures; and the organization bears
the cost of storing sensitive drafts. A serious pilot records each component.
The product earns adoption only if the attention saved at the reader boundary
exceeds those costs for a document class that recurs.

\section{Lessons from the case}
\label{sec:case-lessons}

The worked note establishes five design requirements. The system needs a clear
statement of the audience and decision because purpose is not recoverable from
source syntax alone. It needs located findings because an author cannot assess
an unexplained score. It needs protected comparisons because fluency can conceal
a changed number or qualification. It needs a reader who has not seen the
repository because project familiarity can hide a failure. And it needs an
economic record because a nicer draft is not a product benefit until it reduces
scarce attention without creating a meaning error.

The case does not establish that the workflow improves a real proposal. It uses
a finite teaching scenario and a deliberately small reader protocol. The
comparative program in Part~IV supplies the next observation. The value of the
case is that it makes the proposed object testable: a later study can disagree
about the effect without disagreeing about what the system was supposed to do.

The internal record and the external literature now give the case a wider
test. Each new tool or result can be placed against the same questions: does it
clarify the decision, preserve a protected object, reduce author or reader
burden, or improve the evidence for the next choice? That is the product
boundary in working form.
